\documentclass[11pt]{article}
\usepackage[OT1]{fontenc}
\renewcommand{\rmdefault}{cmr}
\renewcommand{\sfdefault}{cmss}
\renewcommand{\ttdefault}{cmtt}
\usepackage{amsmath,amssymb,amsthm,mathtools}
\usepackage[margin=27mm]{geometry}
\usepackage{microtype,booktabs,longtable,array}
\usepackage[hidelinks]{hyperref}
\setlength{\emergencystretch}{2em}
\setlength{\parskip}{3pt plus 1pt}
\setcounter{tocdepth}{1}
\newtheorem{theorem}{Theorem}[section]
\newtheorem{lemma}[theorem]{Lemma}
\newtheorem{proposition}[theorem]{Proposition}
\newtheorem{corollary}[theorem]{Corollary}
\theoremstyle{definition}
\newtheorem{definition}[theorem]{Definition}
\newtheorem{example}[theorem]{Example}
\theoremstyle{remark}
\newtheorem{remark}[theorem]{Remark}
\newcommand{\N}{\mathbb N}
\newcommand{\Q}{\mathbb Q}
\newcommand{\R}{\mathbb R}
\newcommand{\Z}{\mathbb Z}
\newcommand{\supp}{\operatorname{supp}}
\newcommand{\Con}{\operatorname{Con}}
\newcommand{\ord}{\operatorname{ord}}
\newcommand{\dom}{\operatorname{dom}}
\newcommand{\sgn}{\operatorname{sgn}}
\newcommand{\res}{\operatorname{Res}}
\newcommand{\rad}{\operatorname{rad}}
\newcommand{\CK}{\mathsf{CK}}
\newcommand{\KT}{\mathsf{KT}}
\newcommand{\DCB}{\mathsf{DCB}}
\newcommand{\FS}{\mathsf{FS}}
\newcommand{\Ag}{\mathsf{Ag}}
\newcommand{\Know}{\mathsf K}
\newcommand{\Truth}{\mathsf T}
\newcommand{\Prv}{\mathsf{Pr}}
\newcommand{\code}[1]{\ulcorner #1\urcorner}
\newcommand{\Tot}{\operatorname{Tot}}
\numberwithin{equation}{section}

\title{A Theorem Which Is Not Common Knowledge\\in an Untyped Theory of Knowledge}
\author{}
\date{28 October 2026\\Working manuscript}
\begin{document}
\maketitle
\begin{abstract}
Gorbow's theory KT combines an arithmetized knowledge predicate,
compositional truth, veracity and theorem-only necessitation. Common
knowledge is represented by a self-referential fixed point. Gorbow asks
whether the common-belief necessitation theorem for the subsystem DCB
extends to KT. We give a negative answer for the published theory with
base PA in the expanded language. There is a truth-free sentence
$\theta$ such that KT proves $\theta$ and proves the negation of its
common-knowledge assertion, although every fixed finite iteration of
everyone-knows applied to $\theta$ is a KT theorem. More generally, for
each permitted selector, a local-consistency fixed point is equivalent
in DCB to everyone selected knowing a contradiction. In KT it is
equivalent to the selected group being empty. We prove this classification
and supply a direct one-agent revision argument for consistency of KT
and KT with the untyped Barcan principle. All necessitation rules are
used only on theorems.
\end{abstract}
\tableofcontents

\section{The common-knowledge necessitation question}

Ordinary necessitation says that once a sentence is a theorem, one may
infer that an agent knows it. Common-knowledge necessitation would allow
a further inference: every selected agent knows the sentence, knows
that every selected agent knows it, and so on in the sense formalized
by the common-knowledge predicate.

In a typed modal system these principles often fit together naturally.
Gorbow's setting \cite{Gorbow} is deliberately untyped. Knowledge and
truth apply to codes of sentences in the full language, including
sentences which themselves discuss knowledge and truth. Common knowledge
is expressed by a diagonal fixed point inside arithmetic. These choices
permit useful self-reference but make the passage from individual to
common necessitation a substantive theorem, rather than a matter of
notation.

The common-belief subsystem DCB has a theorem-to-common-belief rule.
In his concluding discussion, item~5(d), Gorbow asks how its results
extend to the larger knowledge theory KT, including the corresponding
common-knowledge rule. We show that the unrestricted extension fails.
The counterexample is particularly sharp: the proposed common knowledge
is not merely unproved; its negation is a theorem.

The result concerns the published full-language rule. An arithmetic-only
rule, a different fixed-point constructor, or an additional restriction
on selectors would require a separate analysis.

\section{The source theories and coding conventions}

Work in Gorbow's expanded first-order language $L$. It contains the
language of arithmetic, an agent predicate $\Ag$, a binary knowledge
predicate $\Know(a,u)$, a unary knowledge predicate for all agents,
a truth predicate $\Truth(u)$, and standard definitional syntax coding.
The sublanguage $L_-$ consists of arithmetic, the source's universal
coding relation and function $U,u$, and $\Ag$; it excludes both knowledge
and truth. Its interpretations will remain fixed in the consistency
construction.
The base in our consistency theorem is exactly $\mathsf{PA}(L)$ with
the source's definitional syntax apparatus. In particular induction is
available for formulas in $L$. No arbitrary extra base theory is assumed.

For a literal sentence $\sigma$, $\code\sigma$ abbreviates its numeral
code. Codes, substitution functions and the numeral denoting a code
must be distinguished: $\CK_A(\code\sigma)$ is a sentence and has its
own code $\code{\CK_A(\code\sigma)}$.

Put $D=\DCB$. Its relevant axioms, with the source sentencehood guards
understood, are
\begin{align}
 &\exists a\,\Ag(a),\label{eq:nonempty}\\
 &\Know(a,\code{\sigma\to\tau})\wedge\Know(a,\code\sigma)
       \ \longrightarrow\ \Know(a,\code\tau)
       \quad(\Ag(a)),\label{eq:distribution}\\
 &\Prv_D(\code\sigma)\ \longrightarrow\ \Know(a,\code\sigma)
       \quad(\Ag(a)).\label{eq:reflection}
\end{align}
These are instances of universally coded axioms, not restrictions of
their domains to an externally chosen finite collection. The unary
knowledge predicate is linked by the source bridge to
$\forall a(\Ag(a)\to\Know(a,u))$.

The index of $\Prv_D$ in \eqref{eq:reflection} is $D$ itself.
It is not the provability predicate for KT. Gorbow constructs its
recursive presentation by a syntactic fixed point. We use that
presentation and the ordinary derivability conditions. In particular
L\"ob's theorem applies to $D$: if
\[
 D\vdash\Prv_D(\code\delta)\to\delta,
 \quad\text{then}\quad D\vdash\delta.
\]
This is the standard metatheorem for provability, not an extra
knowledge axiom. Its application does not require a consistency
assumption.

The theory $\FS$ adds uniform compositional truth axioms for the fixed
$L_-$ atoms, Boolean connectives and quantifiers. The atomic clauses
apply to closed arithmetic term instances; the quantified clauses use
closed $L$-term substitution. It also has the two theorem rules
\[
 \frac{\vdash\sigma}{\vdash\Truth(\code\sigma)},
 \qquad
 \frac{\vdash\Truth(\code\sigma)}{\vdash\sigma}.
\]
For example the Boolean clauses include
$\Truth(\code{\neg\sigma})\leftrightarrow\neg\Truth(\code\sigma)$
and
$\Truth(\code{\sigma\wedge\tau})\leftrightarrow
(\Truth(\code\sigma)\wedge\Truth(\code\tau))$.
Their full coded versions, not just these literal instances, are part
of $\FS$.

Finally KT extends $D$ and $\FS$ by veracity and theorem-only knowledge
necessitation:
\begin{equation}\label{eq:veracity}
 \Ag(a)\to(\Know(a,\code\sigma)\to\Truth(\code\sigma)),
 \qquad
 \frac{\vdash\sigma}{\vdash\forall a(\Ag(a)\to\Know(a,\code\sigma))}.
\end{equation}
No rule is licensed inside a derivation under temporary assumptions.
The source's published figures give the full syntactic schemata;
the displayed formulas isolate exactly the instances used below.

\subsection{Selected agents and common knowledge}

Let $A(x)$ be any permitted $L$-formula with $x$ its only free variable.
For a literal sentence $\sigma$ define
\[
 E_A(\sigma)=\forall x\bigl(\Ag(x)\to
                (A(x)\to\Know(x,\code\sigma))\bigr),
 \qquad
 N_A=\exists x(\Ag(x)\wedge A(x)).
\]
The formula $A$ may contain knowledge or truth; no external evaluation
of it is assumed.

We use Gorbow's chosen fixed-point formula $\CK_A$. Its defining
equivalence, after numeral substitution, gives for each literal
$\sigma$ the $D$ theorem
\begin{equation}\label{eq:fixed}
 \CK_A(\code\sigma)\ \leftrightarrow\
 E_A(\sigma)\wedge E_A\bigl(\CK_A(\code\sigma)\bigr).
\end{equation}
The second conjunct therefore involves the code of the actual
common-knowledge sentence. We do not replace it by an informal infinite
conjunction or by a different predicate satisfying an external equation.

\begin{lemma}\label{lem:three}
For fixed literal sentences $\sigma,\tau$:
\begin{enumerate}
\item if $D\vdash\sigma$, then $D\vdash E_A(\sigma)$;
\item $D\vdash E_A(\sigma\to\tau)\to(E_A(\sigma)\to E_A(\tau))$;
\item $D\vdash\Prv_D(\code\sigma)\to E_A(\sigma)$.
\end{enumerate}
\end{lemma}
\begin{proof}
For the first assertion, a fixed actual proof of $\sigma$ has its
provability certified by the arithmetic base. Apply
\eqref{eq:reflection} and restrict to selected agents. The second
assertion is pointwise \eqref{eq:distribution} followed by first-order
logic. The third is \eqref{eq:reflection} with the same restriction.
No distribution through a quantified knowledge claim is used.
\end{proof}

\section{The local-consistency fixed point collapses}

Let $b$ be the arithmetic contradiction $0=S0$, and set
\[
 \theta_A=(E_A(b)\to b),\qquad c_A=\CK_A(\code{\theta_A}).
\]
The sentence $\theta_A$ says that if all selected agents know the
contradiction, the contradiction holds. It has the form of a local
consistency assertion for the selected knowledge predicates.

\begin{theorem}\label{thm:collapse}
For every allowed selector $A$,
\[
 D\vdash c_A\leftrightarrow E_A(b).
\]
\end{theorem}
\begin{proof}
Write $\delta=(c_A\to E_A(b))$. Reason in ordinary first-order
deduction over $D$ under temporary assumptions
$\Prv_D(\code\delta)$ and $c_A$. Lemma~\ref{lem:three}(3) yields
$E_A(\delta)$, while \eqref{eq:fixed} gives
$E_A(c_A)$ and $E_A(\theta_A)$. Two applications of
Lemma~\ref{lem:three}(2) give, successively,
\[
 E_A(E_A(b)),\qquad E_A(b).
\]
The first uses $\delta=(c_A\to E_A(b))$; the second uses
$\theta_A=(E_A(b)\to b)$. Discharging the assumptions proves
$D\vdash\Prv_D(\code\delta)\to\delta$. L\"ob's theorem now gives
$D\vdash\delta$, proving the forward implication. No necessitation
rule has been applied to an assumption.

For the converse, $D$ proves the two tautological implications
$b\to\theta_A$ and $b\to c_A$. Lemma~\ref{lem:three}(1),(2) gives
\[
 E_A(b)\to E_A(\theta_A),\qquad E_A(b)\to E_A(c_A).
\]
Their conjunction and \eqref{eq:fixed} imply $E_A(b)\to c_A$.
\end{proof}

\begin{corollary}
The same fixed point is $D$-equivalent to common belief of the
contradiction: $D\vdash c_A\leftrightarrow\CK_A(\code b)$.
\end{corollary}
\begin{proof}
The defining equation for $\CK_A(\code b)$ gives its implication to
$E_A(b)$. Conversely the $D$ theorem $b\to\CK_A(\code b)$, placed
under $E_A$, gives $E_A(b)\to E_A(\CK_A(\code b))$.
Use the defining equation again, then Theorem~\ref{thm:collapse}.
\end{proof}

\section{Veracity gives the exact selector classification}

\begin{theorem}\label{thm:selector}
In the original KT theory,
\[
 \KT\vdash\theta_A\leftrightarrow N_A,
 \qquad
 \KT\vdash c_A\leftrightarrow\neg N_A.
\]
\end{theorem}
\begin{proof}
The atomic truth clause for the closed equality $b$ gives
$\Truth(\code b)\leftrightarrow b$. Since PA proves $\neg b$,
veracity proves
$\forall x(\Ag(x)\to\neg\Know(x,\code b))$.
First-order logic therefore gives
$E_A(b)\leftrightarrow\neg N_A$. This is valid even when $A$
contains knowledge or truth: no assumption about its extension is
needed. With $\neg b$, the implication $E_A(b)\to b$ is equivalent
to $\neg E_A(b)$, hence to $N_A$. The second equivalence follows from
Theorem~\ref{thm:collapse}.
\end{proof}

Take $A(x)$ to be $x=x$, denoted $\top$. Nonemptiness of agents gives
the explicit truth-free counterexample
\begin{equation}\label{eq:counterexample}
 \theta=\bigl(\forall x(\Ag(x)\to\Know(x,\code b))\to b\bigr),
 \qquad
 \KT\vdash\theta,\quad
 \KT\vdash\neg\CK_\top(\code\theta).
\end{equation}
Thus adding unrestricted common-knowledge necessitation to KT makes
it inconsistent. Once consistency of KT is established, the displayed
negation also yields nonprovability of the common-knowledge assertion
in KT itself.

\begin{corollary}\label{cor:finite}
For every fixed external $r\in\N$, $\KT\vdash E_\top^r(\theta)$,
where $E_\top^0(\theta)=\theta$ and
$E_\top^{r+1}(\theta)=E_\top(E_\top^r(\theta))$.
\end{corollary}
\begin{proof}
Starting from \eqref{eq:counterexample}, apply theorem-only knowledge
necessitation $r$ times. Each premise is an already established theorem.
\end{proof}

This is a separate theorem for every standard finite $r$. It does not
assert a formula internally quantifying over iteration lengths. In
particular it supplies no inference from all finite iterations to the
chosen self-referential common-knowledge sentence. That missing
inference is exactly what the example exposes.

The empty-selector case is a useful check. If KT proves $\neg N_A$,
it proves $c_A$ and $\neg\theta_A$. If the nonemptiness of the selected
group is undecided, the two equivalences of Theorem~\ref{thm:selector},
rather than an externally chosen case, are the result.

\section{A direct consistency argument}

Gorbow proves consistency by generalized revision semantics
\cite[Theorem~15 and Corollary~16]{Gorbow}. We include a direct
one-agent specialization, both to make the nonprovability conclusion
available here and to explain why the preceding results coexist.
The construction uses standard external satisfaction. It does not
provide one standard structure satisfying every KT theorem.

\begin{theorem}\label{thm:consistency}
The source theories KT and $\KT+\mathsf{UBF}$ are consistent when
their base is $\mathsf{PA}(L)$ with its standard definitional syntax
apparatus. Here $\mathsf{UBF}$ is the source untyped Barcan principle
for universally known closed-arithmetic-term instances.
\end{theorem}
\begin{proof}
Fix an $L_-$ expansion $M$ of the standard natural numbers satisfying
the defining arithmetic syntax axioms. Interpret $\Ag$ by $\{0\}$,
and keep all functions, syntax devices and $L_-$ relations fixed.
Every instance of full $L$-induction holds in every expansion of
the standard natural numbers, whatever interpretations are later given
to knowledge and truth.

Let $P_0$ be the set of codes of literal $D$-theorems. This set exists
without an initial consistency assumption. Define $M_n$ to have truth
extension $P_n$, binary knowledge extension $\{(0,u):u\in P_n\}$,
and unary all-agents knowledge extension $P_n$. Recursively set
\[
 P_{n+1}=\{\code\sigma:M_n\models\sigma\},
\]
where $\sigma$ ranges over actual $L$-sentences. At each standard
finite stage,
\begin{equation}\label{eq:revision}
 M_{n+1}\models\Truth(\code\sigma)
 \quad\Longleftrightarrow\quad M_n\models\sigma.
\end{equation}
The same equivalence holds for $\Know(0,\code\sigma)$ and for unary
knowledge. Satisfaction on the right is external classical satisfaction.

\emph{DCB holds at every stage.} At stage zero the theorem set is
closed under modus ponens, so knowledge distribution holds. The
provability predicate has its standard interpretation, so
\eqref{eq:reflection} holds. Nonemptiness and the unary/binary bridge
hold by construction; the arithmetic base holds as just explained.
Thus $M_0\models D$. If $M_n\models D$, ordinary first-order
soundness says every $D$-theorem is true in $M_n$, hence its code is
in $P_{n+1}$. This gives \eqref{eq:reflection} at the next stage.
Classical truth in $M_n$ is closed under modus ponens, giving
distribution there too. Induction proves $M_n\models D$ for all $n$.
The soundness step concerns the ordinary first-order axiom theory $D$;
it does not insert KT's additional theorem rules into $D$.

\emph{Truth composition holds at positive stages.} For $n\ge1$,
membership in $P_n$ means truth in $M_{n-1}$. Classical satisfaction
immediately gives all Boolean truth clauses in $M_n$, including for
sentences containing knowledge or truth. The $L_-$ atomic clauses
hold because their relations and closed arithmetic term values remain
fixed. For quantifiers, the substitution lemma applies to the fixed
function interpretations. Closed $L$-terms denote natural numbers,
and every natural number has a numeral. Thus truth of all closed-term
instances in $M_{n-1}$ is exactly universal satisfaction there; the
existential clause is similar. This verifies the full coded quantified
clauses, since the ambient arithmetic universe is standard. Veracity
holds even at stage zero, because the single agent's knowledge extension
equals the truth extension. It does not assert that $P_n$ is the set
of sentences true in $M_n$.

Call a sentence \emph{eventually valid} when it is true in all $M_n$
beyond some standard finite stage depending on the sentence. Every
$D$ axiom has threshold zero; every FS axiom has threshold at most
one; veracity has threshold zero. Logical consequences of finitely
many eventually valid sentences are eventually valid, by ordinary
soundness in each sufficiently late structure.

If $\sigma$ is eventually valid, \eqref{eq:revision} makes
$\Truth(\code\sigma)$ eventually valid, and also gives eventual
validity of $\forall a(\Ag(a)\to\Know(a,\code\sigma))$.
Conversely eventual validity of $\Truth(\code\sigma)$ implies
eventual validity of $\sigma$ by shifting the stage back one.
All three theorem rules preserve eventual validity. Induction over
an actual finite KT proof therefore makes its conclusion eventually
valid. The arithmetic falsehood $b$ is false at every stage, so it
cannot be a theorem.

For the additional Barcan axiom, take $n\ge1$. If the sole agent
knows every closed arithmetic term instance of $\varphi(v)$ at stage
$n$, then $M_{n-1}$ satisfies all these instances. Numerals cover the
standard universe, so $M_{n-1}\models\forall v\varphi(v)$.
Equation~\eqref{eq:revision} says the agent knows that universal
sentence at stage $n$. Thus the actual full-language UBF schema is
eventually valid. The same proof induction establishes consistency
of the extension separately.
\end{proof}

The proof uses no internally defined truth iteration, no common
knowledge converse principle, and no assertion that finite revision
stages stabilize globally. Its conclusion is external consistency,
not an internal proof of $\Con(\KT)$.

\section{What the negative answer establishes}

Combining Theorem~\ref{thm:consistency} with
\eqref{eq:counterexample} gives
\[
 \KT\nvdash\CK_\top(\code\theta),
 \qquad
 \KT+\mathsf{UBF}\nvdash\CK_\top(\code\theta).
\]
Both theories nevertheless prove $\theta$ and the negation of that
common-knowledge sentence. This answers the unrestricted
common-knowledge necessitation question in the specified source case.

The selector theorem is slightly more informative than one
counterexample: it identifies the exact role of nonemptiness. The
DCB collapse also explains the obstruction structurally. A fixed point
claiming common belief of local consistency allows provability in DCB
to trigger L\"ob's theorem; it collapses to common belief of the
contradiction. Adding veracity then turns that condition into emptiness
of the selected group.

Strengthening the base by arbitrary axioms preserves the displayed
syntactic derivations, but the consistency proof does not automatically
extend. Likewise, alternative restricted necessitation rules are not
classified by this counterexample. These boundaries keep the result
aligned with the full-language question actually answered.

\paragraph{Preparation.}
This AI-assisted hand-proof manuscript includes the local fixed-point,
selector and consistency arguments. The source theories and L\"ob's
metatheorem are credited mathematical inputs. No formal certification
or established priority is claimed.

\begin{thebibliography}{9}
\bibitem{Gorbow} P.~Gorbow,
\emph{A genuinely untyped explanation of common belief and knowledge},
The Review of Symbolic Logic 18 (2025), no.~1, 160--185.
\href{https://doi.org/10.1017/S175502032400025X}{doi:10.1017/S175502032400025X}.
We use the published DCB/FS/KT axioms, common-belief fixed point,
Theorem~2, Theorem~15, Corollary~16 and concluding discussion~5(d).
\bibitem{Lob} M.~H.~L\"ob,
\emph{Solution of a problem of Leon Henkin},
Journal of Symbolic Logic 20 (1955), no.~2, 115--118.
\end{thebibliography}
\end{document}
